% \iffalse meta-comment
%
% Copyright (C) 2020-2021 by Antoine Missier <antoine.missier@ac-toulouse.fr>
% ----------------------------------------------------------------------
% 
% This file may be distributed and/or modified under the
% conditions of the LaTeX Project Public License, either version 1.3
% of this license or (at your option) any later version.
% The latest version of this license is in:
%
%    http://www.latex-project.org/lppl.txt
%
% and version 1.3 or later is part of all distributions of LaTeX 
% version 2005/12/01 or later.
%
% \fi
%
% \iffalse
%<*driver>
\ProvidesFile{exesheet.dtx}
%</driver>
%<class|package>\NeedsTeXFormat{LaTeX2e}[2005/12/01]
%<class>\ProvidesClass{exesheet}
%<package>\ProvidesPackage{exesheet}
%<*class|package>
[2021/11/15 v1.3 .dtx exesheet file]
%</class|package>
%
%<*driver>
\documentclass{ltxdoc}
% uncomment the following lines if you want to use the exesheet class instead of ltxdoc
%\usepackage{doc}
%\AtBeginDocument{\MakeShortVerb{\|}}
%\providecommand\marg[1]{{\ttfamily\{}\meta{#1}{\ttfamily\}}}
%\providecommand\oarg[1]{{\ttfamily[}\meta{#1}{\ttfamily]}}
\usepackage[notoc]{exesheet}
\usepackage{schooldocs}
\usepackage[utf8]{inputenc}
\usepackage[T1]{fontenc}
\usepackage[english]{babel}
\usepackage{lmodern}
\usepackage{multicol}
\usepackage{hyperref}
\usepackage{sectsty}
\usepackage{pstricks}
\psset{unit=1cm}
\geometry{textwidth=355pt,vmarginratio=1:1,hmarginratio=3:2}
% altered by schoodocs, now very close from the default geometry used by ltxdoc

\newcommand{\ds}{\displaystyle}
\newcommand{\bs}{\textbackslash}

\definecolor{darkbrown}{rgb}{0.5,0.1,0.1}
\allsectionsfont{\sffamily\color{darkbrown}}
\renewcommand{\labelexercisestyle}{\rmfamily\color{black}}
\renewcommand{\labelsubpartstyle}{\rmfamily\color{black}}
\renewcommand{\annexstyle}{\rmfamily\color{black}\MakeUppercase}
\renewcommand{\correctionstyle}{\rmfamily\color{correctioncolor}}

\displaypts
%\diplaynotes gives a non desired margin ratio here
% following instructions simulate \displaynotes without altering the margin ratio
\setlength{\marginparwidth}{\oddsidemargin}
\addtolength{\marginparwidth}{1in}
\addtolength{\marginparwidth}{-\marginparsep}
\setlength{\marginparwidth}{0.8\marginparwidth}
\setboolean{marginpoints}{true}
\setboolean{marginfullnotes}{true}

\DisableCrossrefs
%\CodelineIndex
%\RecordChanges

\title{The \textsf{exesheet} class and package}%\thanks{This document
%    corresponds to \textsf{exesheet}~\fileversion, dated \filedate.}}
%\author{Antoine Missier \\ \texttt{antoine.missier@ac-toulouse.fr}}
\subject{Antoine Missier \\
    \texttt{antoine.missier@ac-toulouse.fr}}
\date{\filedate, \fileversion}
\hypersetup{%
    colorlinks,
    linkcolor=blue,
    pdftitle={exesheet},
    pdfsubject={LaTeX class and package},
    pdfauthor={Antoine Missier}
}
\begin{document}
\DocInput{exesheet.dtx}
\end{document}
%</driver>
% \fi
%
% \CheckSum{613}
%
% \CharacterTable
%  {Upper-case    \A\B\C\D\E\F\G\H\I\J\K\L\M\N\O\P\Q\R\S\T\U\V\W\X\Y\Z
%   Lower-case    \a\b\c\d\e\f\g\h\i\j\k\l\m\n\o\p\q\r\s\t\u\v\w\x\y\z
%   Digits        \0\1\2\3\4\5\6\7\8\9
%   Exclamation   \!     Double quote  \"     Hash (number) \#
%   Dollar        \$     Percent       \%     Ampersand     \&
%   Acute accent  \'     Left paren    \(     Right paren   \)
%   Asterisk      \*     Plus          \+     Comma         \,
%   Minus         \-     Point         \.     Solidus       \/
%   Colon         \:     Semicolon     \;     Less than     \<
%   Equals        \=     Greater than  \>     Question mark \?
%   Commercial at \@     Left bracket  \[     Backslash     \\
%   Right bracket \]     Circumflex    \^     Underscore    \_
%   Grave accent  \`     Left brace    \{     Vertical bar  \|
%   Right brace   \}     Tilde         \~}
%
%
% \changes{v0.0}{2010/09/01}{Initial personal version}
% \changes{v0.9}{2020/01/02}{Initial documentation in french}
% \changes{v1.0}{2020/07/22}{Initial public version with documentation in english}
% \changes{v1.1}{2021/08/30}{Macro tablenum replaced by tablenum1, small modifications due
% to changes in tasks package}
% \changes{v1.2}{2021/09/18}{Bug correction in macros tablenum1 and tablenuma, using the translations package for internationalization}
%\changes{v1.3}{2021/11/15}{Bad spacing in the exe macro and bug correction in the starred version of this macro}
%
% \GetFileInfo{exesheet.dtx}
%
% \maketitle
% \tableofcontents
%
% \section{Introduction}
%
% The \textsf{exesheet} \emph{package} is used for typesetting exercise or exam sheets.
% In addition, the \textsf{exesheet} \emph{class} loads the \textsf{schooldocs} package.
% This one makes adjustments for margins and title and defines various layout styles 
% with particular header and footer, appropriate for exercise sheets (among others).
% See the \textsf{schooldocs} documentation for more details.
% The \textsf{exesheet} \emph{class} is based on \textsf{article} 
% and has the same options.
%
% Many other packages are dedicated to exercise sheets.
% Most propose to encapsulate each exercise in an environment
% while \textsf{exesheet} begins each exercise with the |\exercise| command, 
% which works like a subsection (with same features)
% and is suitable for documents consisting exclusively of exercises.
% The package provides also alternative formatting, 
% more relevant for short exercises.
%
% Another specificity of the \textsf{exesheet} package is specific settings
% for enumeration lists, different from \LaTeX\ standard settings,
% useful for the numbering of questions or answers inside an exercise.
%
% Other packages provide often more or less elaborate mechanisms
% to manage the placement of answers.
% \textsf{exesheet} has no such ambitions:
% for all exercises of the sheet, we can display subject only, 
% answer only or both, but always at the place they are inserted
% in the source file. On the other hand this choice may be very flexible:
% we can do a correct version for all exercises together, or a correct version per exercise,
% per part (subpart of exercise), per question, per sub-question.
%
% Finally the original functionality of this package is the display 
% of a detailed marking scheme,
% with optional explanations, in the margin of exercise answers.
%
% There are two possible options: \texttt{notoc} (see \ref{par:notoc})
% and \texttt{nosetlist} (see \ref{par:nosetlist}).
%
%\medskip
% In the current document, a frame is used to highlight 
% examples beside the rest of the documentation.
%
% \section{Titles}
% \subsection{The \texttt{\bs exercise} command}

% \DescribeMacro{\exercise}
% Each exercise begins with the |\exercise|\oarg{opt} command. 
% This command typesets \textbf{\exercisename}, as a document
% subsection, followed by automatic numbering,
% unique for the whole document.
% The optional parameter \meta{opt} is used to put additional text
% on the same title line, for example to precise a subject or a marking scheme.
% |\exercise[ (to begin)]| yields:
%
% \exercise[ (to begin)]
% \psframe(-0.25,1.3)(12.75,-0.4)
% Try to use this first command now, it's easy.
%
% \bigskip
% |\hrulefill| can be put in the optional argument to produce an horizontal rule.
%
% To bring additional text closer to the exercise number, 
% we can use the |\unskip| command
% which eliminates preceding space, see the following example, obtained with
% |\exercise[\unskip : calculation\ \hrulefill]|:
% \exercise[\unskip: calculation \ \hrulefill]
% \psframe(-0.25,1.3)(12.75,-0.4)
% Calculate $1+1$.
%
% \bigskip
% \DescribeMacro{\exercisename}
% The word \emph{\exercisename} has automatic translation 
% in a few languages\footnote{Translation 
% is currently integrated into the package for the following languages: French,
% German, Spanish, Italian, Portuges.}
% according to the loaded language (by \textsf{babel} or \textsf{polyglossia}).
% It can be redefined, with |\renewcommand|,
% or else we can use macros from the \textsf{translations} 
% package (allowing dynamic language change), e.g. (in the preamble)\\
% |\DeclareTranslation{Swedish}{exesheet-exercise}{\"Ovning}}|.
%
% \DescribeMacro{\labelexercise}
% This command calls |\exercisename| following by the exercise number.
% It can be redefined. For example, to add a period after the exercise number:
% |\renewcommand{\labelexercise}{\exercisename~\theexercise.}|
%
% \DescribeMacro{\theexercise}
% To change only the numbering type, redefine
% the |\theexercise| command, based on the |exercise| counter.
%
% \DescribeMacro{\labelexercisestyle}
% This macro (which is empty by default) allows to define a particular style
% for exercise titles. In the present document, we defined in the preamble:
% |\renewcommand{\labelexercisestyle}{\rmfamily\color{black}}|\footnote{In the 
% present document, to highlight real sections and subsections titles,
% their color and font have been modified with the
% \texttt{\bs allsectionsfont} macro from the \textsf{sectsty} package.}.
%
% \bigskip
% \DescribeMacro{\exercise*}
% The starred version |\exercise*|\oarg{opt}\marg{label}
% allows to choose another \meta{label} for a particular exercise
% and removes the numbering. For instance:
% |\exercise*[(Fermat theorem)]{Problem}| yields:
%
% \exercise*[(Fermat theorem)]{Problem}
% \psframe(-0.25,1.3)(12.75,-0.7)
% Prove that there are no positive integers $x, y, z$ such that $x^n+y^n=z^n$
% for any integer $n$ greater than 2.
%
% \subsection{The \texttt{\bs subpart} command}
%
% \DescribeMacro{\subpart}
% An exercise may contain several parts that we obtain 
% with the |\subpart|\oarg{opt} command, typeset like a sub-subsection.
% \exercise
% \subpart[(preliminary)]
% \psframe(-0.25,1.95)(12.75,-1.9)
% First of all, prepare your cup of tea.
% \subpart
% Now you are ready to make the current exercise.
%
% \bigskip
% \DescribeMacro{\thesubpart}
% By default, the subpart numbering uses letters : A, B, C, etc.
% This numbering type can be redefined with the |\thesubpart| command based on
% the |subpart| counter, for instance |\renewcommand\thesubpart{\arabic{subpart}}|.
%
% \pagebreak
% \DescribeMacro{\subpartname} \DescribeMacro{\labelsubpart} 
% \DescribeMacro{\labelsubpartstyle}
% As for |\exercise|, the |\subpart| command
% uses |\subpartname| (with automatic translation in a few languages 
% according to the selected language),
% |\labelsubpart| and |\labelsubpartstyle|, which are editable.
%
% \DescribeMacro{\subpart*}
% Like |\exercise*|, the starred version |\subpart*|\oarg{opt}\marg{label}
% allows to freely typeset the subpart \marg{label}, for instance
% |\subpart*{First part}|.
%
% \subsection{The \texttt{\bs annex} command}
% \DescribeMacro{\annex}
% The |\annex|\oarg{opt} command typesets the title \textbf{ANNEX}, 
% in uppercase letters, centered
% and in the subsection style,
% with an optional parameter, added on the same line.
%
% \vspace{-\baselineskip}
% \noindent\psframe(-0.25,-0.3)(12.75,-1.25)
% \annex[(to return)]
%
% \DescribeMacro{\annexname}
% The word \emph{\annexname} has automatic translation in a few languages.
% It can be extended to other languages or modified by redefining |\annexname|
% or with macros from the \textsf{translations} package.
%
% \DescribeMacro{\annexstyle}
% The annex title style is set by the |\annexstyle| macro,
% defined as follows: |\newcommand\annexstyle{\MakeUppercase}|.
% This command may be redefined as one wants.
%
% \subsection{Exercise titles in table of contents 
% (notoc package option)} \label{par:notoc}
% \DescribeMacro{notoc}
% By default, the titles \emph{\exercisename}, \emph{\subpartname} or \emph{\annexname},
% appear in the table of contents
% (or in the pdf file summary when \textsf{hyperref} package is used). 
% To avoid this, you can set the boolean \texttt{notoc} to \texttt{true}
% or call the \textsf{exesheet} package (or class) with \texttt{notoc} option.
% But notice that title optional arguments will always be ignored in the table of contents.
%
% \subsection{Short exercises: the \texttt{\bs exe} command}
% \DescribeMacro{\exe}
% The |\exe| command starts an exercise by the abbreviation \textbf{\exname}.\ 
% followed by the exercise number,
% without using sectioning commands, and the exercise body begins on the same
% line. An exercise starts a new paragraph without indentation.
%
% \exe This is a short exercise who
% can contain several paragraphs or questions however.
%
% \noindent\psframe(-0.25,1.4)(12.75,-1.2)
% \indent Here for example starts a new paragraph.
% \exe This is another short exercise.
%
% \bigskip
% \DescribeMacro{\exname} \DescribeMacro{\exlabel} \DescribeMacro{\exsepmark}
% The abbreviation \emph{\exname} may be changed by
% redefining |\exname| or with macros from the \textsf{translations} package.
% The |\exlabel| macro calls |\exname| following by a period then the exercise number,
% and |\exsepmark| typesets a long dash.
% We may change these features by redefining these commands.
%
% \bigskip
% \DescribeMacro{\exe*}
% The starred version prints no separator as shown below:
%
% \vspace{-\baselineskip}
% \noindent\psframe(-0.25,-0.35)(12.75,-1.1) 
% \exe* Another short exercise without separator.
%
% \subsection{The \texttt{\bs points} command}
% \DescribeMacro{\points}
% The |\points|\marg{pts} command displays the number of points awarded to an exercise.
% It is intended to be entered in the optional argument of the 
% |\exercise| command\footnote{This command which has to be used
% in the optional argument of \texttt{\bs exercise}, as several others presented below, 
% is incompatible with
% the \textsf{memoir} class which redefines section commands.}.
% |\exercise[\points{5}]| yields:
% \exercise[\points{5}]
% \psframe(-0.25,1.3)(12.75,-0.4)
% Try to read this whole document without drinking tea and you get five points.
%
% \bigskip
% \DescribeMacro{\pointsname} \DescribeMacro{\pointname}
% Like before, the word \emph{points} and in the singular \emph{point}
% (if \meta{pts} is inferior to 2)
% get automatic translation in a few languages
% and are modifiable.
% 
% \DescribeMacro{\pointsstyle} \DescribeMacro{pointscolor}
% The style used by the |\points| command may be redefined with |\pointsstyle|. 
% The color is set by |pointscolor| with the |\definecolor| command
% (from the \textsf{xcolor} package by Uwe Kern, loaded by \textsf{exesheet}), 
% for instance: \\ |\definecolor{pointscolor}{named}{blue}|.
%
% 
% \section{Enumerations and lists}
%
% \subsection{List settings (nosetlist package option)} \label{par:nosetlist}
% Enumeration lists are intended to represent questions and sub-questions
% inside exercises. For a good highlight,
% labels are typeset in bold. 
% Moreover, they are left aligned, at the start of the line, 
% without indentation, and the vertical space between items
% is increased compared with \LaTeX\ standard lists.
% These settings are done by the |\setlist| command\footnote{Labels
% may also be changed by the \texttt{\bs labelenumi} and \texttt{\bs labelenumii} commands.}
% from the \textsf{enumitem} package by Javier Bezos.
% Lists with \texttt{itemize} environment are kept in their default configuration\footnote{The
% \texttt{french} option of the \textsf{babel} package
% alters \texttt{itemize} lists behavior and uses long dash
% as labels for each list level.
% This behavior is problematic when mathematics follows the dash
% symbol because the latter may be confused with the minus sign.
% Modifications of the \texttt{french} option (which supersedes \texttt{frenchb} 
% of \textsf{babel}) are canceled to restore \LaTeX\ default settings.}.
%
% \vspace{-\baselineskip}
% \noindent\psframe(-0.25,-0.35)(12.75,-3.9) 
% \exercise
% \begin{enumerate}
% \item First question
% \begin{enumerate}
% \item First sub-question
% \item Second sub-question
% \end{enumerate}
% \item Second question
% \end{enumerate}
%
% \DescribeMacro{nosetlist}
% We can avoid enumeration list modifications and restore the \LaTeX\ 
% default settings by passing the option
% \texttt{nosetlist} to the \texttt{exesheet} class or package.
%
% \subsection{List of exercises : the \texttt{exenumerate} environment}
% \DescribeEnv{exenumerate}
% When an exercise sheet is made of short independent exercises,
% it would be ill-advised to display the complete title \emph{\exercisename} for each. 
% In addition to the |\exe| command, previously presented,
% we provide an even lighter solution with the |exenumerate| environment, 
% which is only an enumeration list in which
% spaces between items are further increased compared with those of \texttt{enumerate}.
% Here is an example (the main list is an \texttt{exenumerate} environment but 
% sub-lists are produced with regular \texttt{enumerate} environments):
%
% \vspace{-\baselineskip}
% \noindent\psframe(-0.25,-0.3)(12.75,-6.75) 
% \begin{exenumerate}
% \item Translate the following sentences in english: 
% \begin{enumerate}
% \item Nam dui ligula, fringilla a, euismod sodales, sollicitudin vel, wisi.
% \item Nam lacus libero, pretium at, lobortis vitae, ultricies et, tellus.
% \end{enumerate}
% \item  Translate the following sentence in german:
% \par Donec aliquet, tortor sed accumsan bibendum, erat ligula aliquet magna,
% vitae ornare odio metus a mi. 
% \item  Translate the following sentences in french: 
% \begin{enumerate}
% \item Cum sociis natoque penatibus 
% et magnis dis parturient montes, nascetur ridiculus mus.
% \item Nulla ullamcorper vestibulum turpis.
% \end{enumerate}
% \end{exenumerate}
% The environment takes an optional parameter, like \texttt{enumerate}, 
% which enables, among others, 
% to typeset alternative list labels, e.g.
% |\begin{exenumerate}[A.]|.
% There are many other options (see the \textsf{enumitem} package documentation).
%
% \subsection{Items aligned by lines}
% \DescribeEnv{tablenum1} 
% \DescribeEnv{tablenuma}
% \DescribeEnv{tablitem}
% These three environments are used to typeset short questions (\texttt{tablenum1}),
% sub-questions (\texttt{tablenuma}) or \texttt{itemize} lists (\texttt{tablitem})
% on the same line. They have the same syntax:
% |\begin{tablenum1}|\oarg{opt}|(|\meta{cols}|)|.
% The \meta{cols} parameter is the number of columns used by the environment. It must be
% \emph{in parentheses}. This parameter can be omitted, then its value is 2.
% As for classic lists, each item begins with the |\item| command.
%
% These three environments are defined by the
% |\NewTasksEnvironment| macro, from the \textsf{tasks} package by Clemens Niederberger.
% They take an optional argument \meta{opt}
% explained in the documentation of this package, 
% e.g.\ |label=\arabic*)| produces an Arabic numbering
% following by a closing parenthesis (like in \texttt{enumitem} package).
% There are also many possibilities to place items in an original way,
% for instance, the |\item*| command allows to
% specify the number of columns the item is supposed to span.
% In the following example,
% the five |\item| are placed in order between
% |\begin{tablenum1}(3)| and |\end{tablenum1}|.
% Notice that numbering is made line by line.
%
% \exercise
% \psframe(-0.25,1.3)(12.75,-2.7)
% Give the derivative of the following functions:
% \begin{tablenum1}(3)
% \item $\ds f(x) = \frac{1-x^2}{\mathrm{e}^x +\mathrm{e}^{-x}}$,
% \item $\ds g(x) = \ln \left( \frac{1-x}{1+x^2}\right)$,
% \item $\ds h(x) = \int_{0}^{1}\mathrm{e}^{xy}\, \mathrm{d}y$,
% \item $\ds k(x) = \sum_{i=1}^{\infty}\frac{1}{x^i}$,
% \item $\ds l(x) = \int_{\frac{1}{x}}^{x}\frac{1}{\ln t}\, \mathrm{d}t$.
% \end{tablenum1}
%
% For \texttt{tablenuma}, labels are
% letters (a, b, c, \ldots) surrounded
% by parentheses. This cannot be modified globally
% (except by redefining the environment with |\RenewTasksEnvironment|).
% If the \textsf{exesheet} package is called with option \texttt{nosetlist},
% labels of \texttt{tablenum1} and \texttt{tablenuma} environments
% are displayed in normal font instead of boldface, with indentation.
%
% \subsection{Items aligned by columns: \texttt{colsenum}}
% \DescribeEnv{colsenum}    
% To get numbering of items by columns, we have the
% \texttt{colsenum} environment: |\begin{colsenum}|\oarg{opt}\marg{cols}.
% The mandatory parameter is the number of columns
% and the optional one will be passed to \texttt{enumerate},
% allowing, for example, to change the type of numbering (a, A, etc.).
% To use this environment, we have to load the
% \textsf{multicol} package in the preamble.
%
% \exercise
% \psframe(-0.25,1.3)(12.75,-2.65)
% Give the derivative of the following functions:
% \begin{colsenum}{3}
% \item $\ds f(x) = \frac{1-x^2}{\mathrm{e}^x +\mathrm{e}^{-x}}$,
% \item $\ds g(x) = \ln \left( \frac{1-x}{1+x^2}\right)$,
% \item $\ds h(x) = \int_{0}^{1}\mathrm{e}^{xy}\, \mathrm{d}y$,
% \item $\ds k(x) = \sum_{i=1}^{\infty}\frac{1}{x^i}$,
% \item $\ds l(x) = \int_{\frac{1}{x}}^{x}\frac{1}{\ln t}\, \mathrm{d}t$.
% \end{colsenum}
% \DescribeEnv{colsenum*}
% We will notice that, on each line, items are not necessarily well
% aligned, what can produce inelegant effects. 
% On the other hand, the \texttt{colsenum} environment doesn't try to align columns 
% from the bottom by adjusting space between items.
% If we want to get this
% (which is the default option in \textsf{multicol}) we have the
% \texttt{colsenum*} environment (same syntax than \texttt{colsenum}).
% Here what we get in that case, with the same exercise.
%
% \exercise
% \psframe(-0.25,1.3)(12.75,-2.65)
% Give the derivative of the following functions:
% \begin{colsenum*}{3}
% \item $\ds f(x) = \frac{1-x^2}{\mathrm{e}^x +\mathrm{e}^{-x}}$,
% \item $\ds g(x) = \ln \left( \frac{1-x}{1+x^2}\right)$,
% \item $\ds h(x) = \int_{0}^{1}\mathrm{e}^{xy}\, \mathrm{d}y$,
% \item $\ds k(x) = \sum_{i=1}^{\infty}\frac{1}{x^i}$,
% \item $\ds l(x) = \int_{\frac{1}{x}}^{x}\frac{1}{\ln t}\, \mathrm{d}t$.
% \end{colsenum*}
% We can see that these alignments are less good than those
% obtained by numbering of items by line. The numbering of items by columns
% may nevertheless be preferable when there are many items with variable heights,
% and a number of items which can be different from one column to the other.
% Moreover, an advantage of \texttt{colsenum} is that the choice of labels 
% is automatic depending on the list level (and the language),
% unlike \texttt{tablenum1} or \texttt{tablenuma}. 
%
% \DescribeEnv{colsitem} 
% \DescribeEnv{colsitem*}
% For \texttt{itemize} lists, the environment \texttt{colsitem}
% produces items aligned by columns rather than by lines as for
% \texttt{tablitem}:
% |\begin{colsitem}|\oarg{opt}\marg{cols}. The optional parameter,
% which is passed to the underlying \texttt{itemize} environment, 
% allows to change the item label (bullet by default).
% And, as for \texttt{colsenum*}, the \texttt{colsitem*} environment
% produces an alignment of columns from the bottom.
%
% \section{Questions and solutions}
%
% \subsection{The \texttt{questions} and \texttt{answers} environments}
% \DescribeEnv{questions} \DescribeEnv{answers}
% The package provides the two environments \texttt{questions} and \texttt{answers}
% to make appear or disappear questions and answers of exercises.
% Display is controlled by two booleans: \texttt{questions} and \texttt{answers}.
% Their default value is \texttt{true}.
%
% \DescribeMacro{\questionsonly} \DescribeMacro{\answersonly}
% The |\questionsonly| command allows to display the questions without the answers
% and |\answersonly| displays the solutions without the questions\footnote{Of course 
% we can also assign the booleans \texttt{questions} and \texttt{answers} 
% directly with the \LaTeX\ command \texttt{\bs setboolean}.}.
% 
% \bigskip
% \DescribeMacro{\correctionstyle} \DescribeMacro{correctioncolor}
% \DescribeMacro{\correctionname}
% In the case (by default) where questions and answers are displayed both,
% answers are then typeset in the style
% |\correctionstyle|, which uses the color \texttt{correctioncolor}.
% This color may be chosen with the |\definecolor| macro
% (by default |\definecolor{correctioncolor}{rgb}{0,0.2,0.6}| = kind of dark blue).
% Moreover, in that case,
% the title \textbf{\correctionname} is displayed at the beginning
% of the \texttt{answers} environment.
% It is defined by the |\correctionname| macro (with translation in a few languages
% and that we can change:\footnote{Except for English, if you want to change package 
% keywords you have to do it \emph{after} \texttt{\bs begin\{document\}}
% because \textsf{babel} will modify these definitions at the end of the preamble.}
% for example we may prefer \emph{Solution}
% than \emph{\correctionname}).
% The style defined by |\correctionstyle| will apply to the title
% but also to the whole environment.
% Here an example.
%
% \vspace{-\baselineskip}
% \noindent\psframe(-0.25,-0.25)(12.75,-6.6) 
% \exercise[\points{3}] \label{points}
% \begin{questions}
% \begin{enumerate}
% \item Is the \textsf{exesheet} package useful ? 
% \item Isn't there any other packages that deal with exercises ?
% \end{enumerate}
% \end{questions}
% \begin{answers}
% \begin{enumerate}
% \item Yes, the \textsf{exesheet} package is very useful.
% \item There are many other packages that deal with exercises, and
% give the ability to produce separately questions and solutions, for example
% \textsf{exercise} by Paul Pichaureau, \textsf{exercises} by Roger Jud,
% \textsf{exsheets} (superseded by \textsf{xsim}) by Clemens Niederberger,
% \textsf{exframe} by Niklas Beisert,
% \textsf{exam} by Philip Hirschhorn,
% \textsf{answers} by Mike Piff and Joseph Wright,
% \textsf{probsoln} by Nicola Talbot, etc.
% \end{enumerate}
% \end{answers}
%
% \DescribeMacro{\points}
% When only answers are displayed, the text color remains black,
% the word \emph{\correctionname} is not displayed and the
% |\points| macro is patched to suppress the display of points.
% An extended solution for a variable display of points will be discussed
% in section \ref{totalpoints} with the |\totalpoints| macro.
%
% \subsection{More about \texttt{answers} environment}
% Internally, we have used the |\comment| and |\endcomment| macros
% from the \textsf{versions} package by Uwe Lück.
% Other excellent packages allow to manage
% selectively piece of code. Let us mention 
% \textsf{verbatim} by Rainer Schöpf, \textsf{comment} by Victor Eijkhout, 
% \textsf{version} by Donald Arseneau and Stephen Bellantoni,
% \textsf{optional} by Donald Arseneau
% and \textsf{codesection} by Matthias Pospiech.
%
% The \textsf{versions} package used provides furthermore 
% the |\includeversion|\marg{env} and
% |\excludeversion|\marg{env} macros who allow to make appear or disappear
% any environment \meta{env} and these optional environments
% may be nested\footnote{The \textsf{codesection} package also allows
% such nesting, including in the preamble,
% as well as the \textsf{optional} package, but the latter only manages short optional code.}.
%
% However the \texttt{questions} and \texttt{answers} environments
% perform another task, not only making appear or disappear piece of text.
% Indeed, there is a problem with the title \emph{\correctionname}. 
% In which format to typeset it and at which level to put it in the table of contents
% (or in the summary of the pdf file)?
% In fact it depends on which level the environment has been nested.
% We can make a single \texttt{answers} environment for all the
% exercises of the sheet or an \texttt{answers} environment
% for each exercise, for each exercise part, for each question or sub-question.
% In fact, the typeset of the title \emph{\correctionname}
% and his level in the table of contents will be adjusted automatically.
%
% \DescribeEnv{\oarg{level}}
% Nevertheless, we can imagine twisted situations in which 
% the title level will not be correctly calculated.
% It is then possible to force the level of the title \emph{\correctionname} 
% with |\begin{answers}|\oarg{level}.
% The optional \meta{level} parameter is defined as follows:
% 1 for section level titles,
% 2 for subsections (like \emph{\exercisename}), 
% 3 for sub-subsections (like \emph{\partname}),
% another number for lower levels
% (then they will not appear in the table of contents).
%
% \bigskip
% \DescribeMacro{answers*}
% The starred version \texttt{answers*} allows to make the title
% \textbf{\correctionname} completely disappear.
%
% \subsection{The \texttt{\bs question} command}
% \DescribeMacro{\question}
% Instead of the \texttt{questions} and \texttt{answers} environments,
% we can also use the simple |\question|\marg{ques}\marg{ans} macro
% in which display of \meta{ques} and \meta{ans} arguments 
% is controlled by the same previous commands |\questionsonly| and |\answersonly|.
% It can be more readable (in the source code) when questions and answers are short.
% In the case of displaying both questions and answers,
% the \emph{\correctionname} title will appear on a new line, in bold,
% as for levels lower than sub-subsection in the \texttt{answers} environment
% (and without entries in the table of contents).
% But this command doesn't work with \texttt{verbatim} inside.
%
% \section{Marginal notes for marking scheme}
%
% The \textsf{exesheet} package allows the display of a marking
% scheme with comments and explanations for answers.
%
% \pagebreak
% \subsection{The \texttt{\bs pts} command}
% \DescribeMacro{\pts}
% When exercises are typeset with the |\exe| macro,
% or as a list with the \texttt{exenumerate} environment,
% the marking scheme is displayed in the margin, on the line where
% we put the |\pts|\marg{num} command (in general the first line of the exercise).
% The \meta{num} parameter is the number of points assigned to the exercise.
%
% \vspace{-\baselineskip}
% \noindent\psframe(-3.5,-0.25)(12.75,-2)
% \exe \pts{3} The first exercise with marking scheme.
% \exe \pts{1.5} The second one.
% 
% \bigskip
% \DescribeMacro{\ptsname} \DescribeMacro{\ptname} 
% \DescribeMacro{ptscolor} \DescribeMacro{\ptsstyle}
% The abbreviation \textsf{\ptsname} (or \textsf{\ptname} 
% when the number of points is inferior to 2)
% is automatically added thanks to the macro |\ptsname| or |\ptname| 
% (translated in a few languages if \textsf{babel} is loaded).
% The display color of the points is defined by \texttt{ptscolor} which can be changed
% with |\definecolor|: by default |\definecolor{ptscolor}{named}{red}|.
% The display style is defined by |\ptsstyle|.
%
% \DescribeMacro{\displaypts}
% The display of the marking scheme, as above, is internally controlled by the
% \texttt{marginpts} boolean.
% The |\displaypts| command, which has to be placed in the preamble, 
% defines \texttt{marginpts} to \texttt{true},
% places marginal notes to the left (and not to the right which 
% is the default behavior of \LaTeX)
% and shifts the text body right to increase the left margin compared 
% with the right one, with a ratio 3:2,
% as it is in the present document.
% This ratio is defined by the
% |geometry| macro from the \textsf{geometry} package (by Hideo Umeki) 
% and is valid for the whole document.
% 
% \subsection{The \texttt{\bs totalexe}, \texttt{\bs note} 
% and \texttt{\bs totalpoints} commands}
% \label{totalpoints}
% \DescribeMacro{\totalexe}
% \DescribeMacro{\note}
% To give a more detailed marking scheme, the \textsf{exesheet} package 
% provides some extended commands:
% |\totalexe|\marg{num} to display the total number of points of an exercise
% and especially |\note|\oarg{pts}\marg{comment},
% to indicate some marking detail in answers.
% The optional parameter \meta{pts} is the number of points of the question
% and the mandatory parameter \meta{comment} is used to give details about the marking scheme.
% Below the title has been obtained with |\exercise[\totalexe{4}]|
% and the first comment (on next page) with \\
% |\note[1]{0,5 for the anti-derivative\\0,5 for simplification of $\ln$}|.
%
% \exercise[\totalexe{4}] \label{totalexe}
% \begin{questions}
% \psframe(-3.5,1.3)(12.75,-3.1)
% For each following question, say if the assertion is true or false.
% Justify the answer carefully.
% \begin{tablenum1}
% \item $\ds \int_{0}^{\sqrt{3}}\frac{1}{x+\sqrt{3}} \,\mathrm{d}x = \ln 2$,
% \item $\ds \int_{2}^{\mathrm{e}} \frac{1}{x\ln x} \,\mathrm{d}x = -\ln 2$,
% \item* The function $F$ defined on $\mathbf{R}$ by 
% $F(x) = \ds \int_{0}^{x}\frac{1}{,t^2+t+1}\,\mathrm{d} t$ is increasing on $\mathbf{R}$.
% \end{tablenum1}
% \end{questions}
% 
% \pagebreak
% \begin{answers}
% \vspace{-\baselineskip}
% \noindent\psframe(-4.7,0.9)(12.75,-10.25)
% \begin{enumerate}
%    \item \note[1]{0,5 for the anti-derivative\\0,5 for simplification of $\ln$}
%    We calculate:
%    \[ \int_{0}^{\sqrt{3}}\frac{1}{x+\sqrt{3}}\,\mathrm{d}x
%        = \bigg[ \ln\left(x+\sqrt{3}\right) \bigg]_0^{\sqrt{3}}
%        = \ln\left(2\sqrt{3}\right)-\ln \sqrt{3}
%        = \ln\left(\frac{2 \sqrt{3}}{\sqrt{3}}\right)
%        =\ln 2.
%    \]
%    \textbf{TRUE}.
%
%    \item \note[1.5]{1 for the anti-derivative \\0.5 for the integral value \\[2ex]
%        we could also see that $\frac{1}{x \ln x} >0$ on $[2, \mathrm{e}]$ 
%       while $-\ln2<0$
%    }
%    We have $\ds\frac{1}{x \ln x}= \frac{\frac{1}{x}}{\ln x}= \frac{u'(x)}{u(x)}$
%    with $u(x)=\ln x$, which is positive on $[2, \mathrm{e}]$. Hence
%
%    \[ \int_{2}^{\mathrm{e}}\frac{1}{x \ln x}\,\mathrm{d}x
%        = \Big[ \ln(\ln x) \Big]_2^\mathrm{e}
%        = \ln(\ln \mathrm{e})-\ln(\ln2)= \ln1 -\ln(\ln2)=-\ln(\ln2).
%    \]
%    \textbf{FALSE}.
%
%    \item \note*{1.5}
%    The function $F$ defined on $\mathbf{R}$ by
%    \[ F(x)= \int_{0}^{x}\frac{1}{t^2+t+1}\, \mathrm{d}t \]
%    is derivable on $\mathbf{R}$ and its derivative is such that
%    $F'(x)=\frac{1}{x^2+x+1}$.
%    \note{0.5 for $F'$\\1 for the sign of $F'$ and conclusion}
%    The denominator is a quadratic polynomial, always positive because its discriminant is
%    $\Delta=-3 <0$.
%    Thus $F$ is increasing on $\mathbf{R}$. \\
%    \textbf{TRUE}.
% \end{enumerate}
% \end{answers}
% In the comment of answer 2, a wider vertical space is produced at line break
% with the optional argument |\\[2ex]|\footnote{When a marginal note at the end of a page,
% is followed by another one at the top of the next page,
% the latter is shifted down.
% We can raise this note, by adding at the line break
% a negative vertical space : \texttt{\bs note*\{\bs\bs[-5ex] 2\}}.}.
% The last comment, not placed beside the number of points of answer 3,
% has been produced by placing on the first line after the formula:
% |\note{0.5 for $F'$\\1 for the sign of $F'$ and conclusion}|.
%
% \DescribeMacro{\note*}
% If we don't want to add comment, we have the
% |\note*|\marg{pts} command (better than |\note|\oarg{pts}|{}|),
% used here in answer 3: |\note*{1.5}|.
%
% \DescribeMacro{markingcolor} \DescribeMacro{\markingstyle} \DescribeMacro{\ptsboxlength}
% The display color of the points, in |\totalexe| and |\note|,
% is defined by \texttt{markingcolor}
% and the style by |\markingstyle|, which are modifiable.
% The oval box produced by |\totalexe| is obtained with the |\ovalbox| command
% of the \textsf{fancybox} package by Timothy Van Zandt,
% with corner arcs set by |\cornersize{1}|.
% The box length is not adjusted to content but depends of the value
% of |\ptsboxlength| in order to keep a uniform appearance
% from one exercise to the other.
%
% \DescribeMacro{notecolor} \DescribeMacro{\notestyle} \DescribeMacro{\noteragged}
% Comment notes are typeset by default in a dark green color
% defined by |\definecolor{notecolor}{rgb}{0.0,0.4,0.0}|.
% The style of the comment is set by the |\notestyle| macro and itself uses the
% |\noteragged| macro. The latter is set by default as |\raggedleft|
% to obtain a \emph{right} alignment (what is customary for notes
% in the left margin).
% 
% \DescribeMacro{\displaypoints}
% Display of points and comments is optional. It is internally controlled by the booleans
% \texttt{marginpoints} and \texttt{marginfullnotes}.
% The |\displaypoints| macro, to place in the preamble, sets  \texttt{marginpoints}
% to \texttt{true}, chooses the left margin for display marginal notes
% and sets the ratio between left and right margins to 3:2\footnote{So 
% that the effect on the margin ratio is correct, the
% \texttt{\bs displaypoints} macro must be placed after other commands 
% that could alter the page geometry, for example the \texttt{\bs pagestyle}
% of the \textsf{schooldocs} package.}.
% Points will then be displayed without the comments.
% 
% \DescribeMacro{\displaynotes} 
% The |\displaynotes|\oarg{align} macro, to be placed also in the preamble,
% defines both booleans to \texttt{true} and allows then the complete display 
% of points and comments.
% The ratio between left/right margins is then set to 5/1\footnote{The
% present document kept a ratio of 3/2 because 
% the \textsf{ltxdoc} class, used here, leaves large enough margins
% and previous examples are not concerned with display of comments.}.
% The optional parameter \meta{align} is used to define alignment:
% |\raggedleft| by default, but we could use |\centering| or
% |\raggedright| (not recommended in left margin).
%
% \DescribeMacro{\displaynotesright}
% The |\displaynotesright|\oarg{align} macro displays a detailed marking scheme
% with comments in the right margin, with the same type of settings 
% as |\displaynotes| but an optional parameter \meta{align} set by default
% to |\raggedleft|.
%
% \bigskip
% \DescribeMacro{\totalpoints}
% The |\totalpoints|\marg{num} macro is intended
% to replace |\points| when using a detailed marking scheme.
% When the scale is not displayed (\texttt{marginpoints} boolean to \texttt{false}),
% it uses |\points| and when the scale is displayed,
% it uses |\totalexe|.
% For example, in the exercise \ref{totalexe}, we should use
% |\totalpoints| rather than |\totalexe|,
% because, when the detailed marking scheme is not displayed, the total points will be typeset
% as in the exercise \ref{points} rather than in the margin.
%
%
%
% \StopEventually{}
%
% \section{Implementation}
%
% \subsection{Class and package basic instructions}
% The \texttt{exesheet} class is based on the \texttt{article} class
% and gives it all its options, except explicit options 
% \texttt{notoc} and \texttt{nosetlist}.
%
%    \begin{macrocode}
%<*class>
\DeclareOption{notoc}{\PassOptionsToPackage{notoc}{exesheet}}
\DeclareOption{nosetlist}{\PassOptionsToPackage{nosetlist}{exesheet}}
\DeclareOption*{\PassOptionsToClass{\CurrentOption}{article}}
\ProcessOptions \relax
\LoadClass{article}
\RequirePackage{exesheet}
\RequirePackage{schooldocs}
%</class>
%    \end{macrocode}
%
% Then we define initialization instructions of the \textsf{exesheet} package.
%
% \medskip
%  The \texttt{shortlabel} option in the \textsf{enumitem} package allows to use labels like in 
% the \textsf{enumerate} package e.g. 1., a), A. etc.
%
%    \begin{macrocode}
%<*package>
\RequirePackage{ifthen}
\newboolean{notoc}
\newboolean{nosetlist}
\DeclareOption{notoc}{\setboolean{notoc}{true}}
\DeclareOption{nosetlist}{\setboolean{nosetlist}{true}}
\ProcessOptions \relax
\RequirePackage{xcolor}
\RequirePackage[shortlabels]{enumitem}
\RequirePackage{tasks}
\RequirePackage{versions}
%\@ifpackageloaded{doc}{}{\RequirePackage{verbatim}}
\RequirePackage{geometry}
\RequirePackage{fancybox}
\RequirePackage{translations}

%    \end{macrocode}
%
% \subsection{Internationalization}
% Here we define the keywords and their translation in French, German, Spanish
% Italian, Portuguese, thanks to macros of the \textsf{translations} package 
% by Clemens Niederberger. It detects the used language 
% loaded by \textsf{babel} or \textsf{polyglossia}.
%
% Accented characters cannot be used here because they are not recognized 
% if \textsf{inputenc} is loaded after \textsf{exesheet}.
% So we have used the basic \LaTeX\ commands to produce them.
%
%    \begin{macrocode}
\DeclareTranslationFallback{exesheet-exercise}{Exercise}
\DeclareTranslationFallback{exesheet-subpart}{Part}
\DeclareTranslationFallback{exesheet-annex}{Annex}
\DeclareTranslationFallback{exesheet-ex}{Ex}
\DeclareTranslationFallback{exesheet-points}{points}
\DeclareTranslationFallback{exesheet-point}{point}
\DeclareTranslationFallback{exesheet-correction}{Correction}
\DeclareTranslationFallback{exesheet-pts}{pts}
\DeclareTranslationFallback{exesheet-pt}{pt}

\DeclareTranslation{English}{exesheet-exercise}{Exercise}
\DeclareTranslation{English}{exesheet-subpart}{Part}
\DeclareTranslation{English}{exesheet-annex}{Annex}
\DeclareTranslation{English}{exesheet-ex}{Ex}
\DeclareTranslation{English}{exesheet-points}{points}
\DeclareTranslation{English}{exesheet-point}{point}
\DeclareTranslation{English}{exesheet-correction}{Correction}
\DeclareTranslation{English}{exesheet-pts}{pts}
\DeclareTranslation{English}{exesheet-pt}{pt}

\DeclareTranslation{French}{exesheet-exercise}{Exercice}
\DeclareTranslation{French}{exesheet-subpart}{Partie}
\DeclareTranslation{French}{exesheet-annex}{Annexe}
\DeclareTranslation{French}{exesheet-ex}{Ex}
\DeclareTranslation{French}{exesheet-points}{points}
\DeclareTranslation{French}{exesheet-point}{point}
\DeclareTranslation{French}{exesheet-correction}{Correction}
\DeclareTranslation{French}{exesheet-pts}{pts}
\DeclareTranslation{French}{exesheet-pt}{pt}

\DeclareTranslation{German}{exesheet-exercise}{\"Ubung}
\DeclareTranslation{German}{exesheet-subpart}{Teil}
\DeclareTranslation{German}{exesheet-annex}{Anhang}
\DeclareTranslation{German}{exesheet-ex}{\"Ub}
\DeclareTranslation{German}{exesheet-points}{Punkte}
\DeclareTranslation{German}{exesheet-point}{Punkt}
\DeclareTranslation{German}{exesheet-correction}{Verbesserung}
\DeclareTranslation{German}{exesheet-pts}{Pkte}
\DeclareTranslation{German}{exesheet-pt}{Pkt}

\DeclareTranslation{Spanish}{exesheet-exercise}{Ejercicio}
\DeclareTranslation{Spanish}{exesheet-subpart}{Parte}
\DeclareTranslation{Spanish}{exesheet-annex}{Anexo}
\DeclareTranslation{Spanish}{exesheet-ex}{Ej}
\DeclareTranslation{Spanish}{exesheet-points}{puntos}
\DeclareTranslation{Spanish}{exesheet-point}{punto}
\DeclareTranslation{Spanish}{exesheet-correction}{Correcci\'on}
\DeclareTranslation{Spanish}{exesheet-pts}{ptos}
\DeclareTranslation{Spanish}{exesheet-pt}{pto}

\DeclareTranslation{Italian}{exesheet-exercise}{Esercizio}
\DeclareTranslation{Italian}{exesheet-subpart}{Parte}
\DeclareTranslation{Italian}{exesheet-annex}{Annesso}
\DeclareTranslation{Italian}{exesheet-ex}{Es}
\DeclareTranslation{Italian}{exesheet-points}{punti}
\DeclareTranslation{Italian}{exesheet-point}{punto}
\DeclareTranslation{Italian}{exesheet-correction}{correzione}
\DeclareTranslation{Italian}{exesheet-pts}{pti}
\DeclareTranslation{Italian}{exesheet-pt}{pt}

\DeclareTranslation{Portuges}{exesheet-exercise}{Exerc\'icio}
\DeclareTranslation{Portuges}{exesheet-subpart}{Parte}
\DeclareTranslation{Portuges}{exesheet-annex}{Anexo}
\DeclareTranslation{Portuges}{exesheet-ex}{Ex}
\DeclareTranslation{Portuges}{exesheet-points}{pontos}
\DeclareTranslation{Portuges}{exesheet-point}{ponto}
\DeclareTranslation{Portuges}{exesheet-correction}{Corre\c c\~ao}
\DeclareTranslation{Portuges}{exesheet-pts}{pts}
\DeclareTranslation{Portuges}{exesheet-pt}{pt}

\newcommand*\exercisename{\GetTranslation{exesheet-exercise}}
\newcommand*\subpartname{\GetTranslation{exesheet-subpart}}
\newcommand*\annexname{\GetTranslation{exesheet-annex}}
\newcommand*\exname{\GetTranslation{exesheet-ex}}
\newcommand*\pointsname{\GetTranslation{exesheet-points}}
\newcommand*\pointname{\GetTranslation{exesheet-point}}
\newcommand*\correctionname{\GetTranslation{exesheet-correction}}
\newcommand*\ptsname{\GetTranslation{exesheet-pts}}
\newcommand*\ptname{\GetTranslation{exesheet-pt}}

%    \end{macrocode}
%
% \subsection{Colors}
% The color \texttt{pointscolor} is used by the macro |\points| 
% and \texttt{markingcolor} by the macros |\totalexe|, |\note*| 
% and also for the optional argument of |\note|;
% \texttt{notecolor} is used by the mandatory argument of |\note|.
%
%    \begin{macrocode}
\definecolor{pointscolor}{named}{red}
\definecolor{ptscolor}{named}{red}
\definecolor{markingcolor}{named}{red}
\definecolor{notecolor}{rgb}{0.0, 0.4, 0.0} % kind of dark green
\definecolor{correctioncolor}{rgb}{0,0.2,0.6} % kind of dark blue

%    \end{macrocode}
%
% \medskip
% \subsection{Titles}
% The \texttt{exercise} counter numbers exercises for the whole document 
% regardless of any section.
% To reset the counter at some point, just write |\setcounter{exercise}{0}|
% and if we want an automatic reset at each section,
% add in the preamble \\
% |\makeatletter \@addtoreset{exercise}{section} \makeatother|.
%
% The parts counter depends on \texttt{exercise} and is reset
% at each new exercise.
%
% The |\labelexercisestyle| and |\labelsubpartstyle| commands
% are empty, but allow to customize the style, for instance:\\
% |\renewcommand\labelexercisestyle{\sffamily}|.
%
% By default, the table of contents displays the titles of exercises and parts.
% The \texttt{notoc} boolean is used to modify this behavior.
% If we want to display only exercise titles but not parts,
% we can place in the preamble |\setcounter{tocdepth}{2}|.
%
% \begin{macro}{\exercise}
%    \begin{macrocode}
\newcounter{exercise}

\newcommand{\labelexercise}{\exercisename~\theexercise}
\newcommand{\labelexercisestyle}{}
\newcommand*{\@exercise}[1][]{%     
    \refstepcounter{exercise}
    \subsection*{\labelexercisestyle\labelexercise\ #1}
    \ifthenelse{\boolean{notoc}}{}{
        \addcontentsline{toc}{subsection}{\labelexercise}
    }
}
\newcommand*{\@@exercise}[2][]{%
    \subsection*{\labelexercisestyle #2 #1}
    \setcounter{subpart}{0} % resets the parts counter
    \ifthenelse{\boolean{notoc}}{}{
        \addcontentsline{toc}{subsection}{#2}
    }
}
\newcommand{\exercise}{\@ifstar{\@@exercise}{\@exercise}}

%    \end{macrocode}
% \end{macro}
% \begin{macro}{\subpart}
%    \begin{macrocode}
\newcounter{subpart}[exercise] %
\renewcommand{\thesubpart}{\Alph{subpart}}

\newcommand{\labelsubpart}{\subpartname~\thesubpart}
\newcommand{\labelsubpartstyle}{}
\newcommand*{\@subpart}[1][]{%
    \refstepcounter{subpart}%    
    \subsubsection*{\labelsubpartstyle\labelsubpart\ #1}
    \ifthenelse{\boolean{notoc}}{}{
        \addcontentsline{toc}{subsubsection}{\labelsubpart}
    }
}
\newcommand*{\@@subpart}[2][]{%
    \subsubsection*{\labelsubpartstyle #2 #1}
    \ifthenelse{\boolean{notoc}}{}{
        \addcontentsline{toc}{subsubsection}{#2}
    }
}
\newcommand{\subpart}{\@ifstar{\@@subpart}{\@subpart}}

%    \end{macrocode}
% \end{macro}
% \begin{macro}{\annex}
%    \begin{macrocode}
\newcommand{\annexstyle}{\MakeUppercase}
\newcommand*{\annex}[1][]{%
    \subsection*{\mbox{}\hfill\annexstyle{\annexname} #1\hfill\mbox{}}
    \ifthenelse{\boolean{notoc}}{}{
        \addcontentsline{toc}{subsection}{\annexname}
    }
}

%    \end{macrocode}
% \end{macro}
% \begin{macro}{\exe}
%    \begin{macrocode}
\newcommand{\exlabel}{\exname.~\theexercise}
\newcommand{\exsepmark}{---}
\newcommand{\@exe}{\bigskip\refstepcounter{exercise}
    \par\noindent\textbf{\exlabel~\exsepmark}~}
\newcommand{\@@exe}{\bigskip\refstepcounter{exercise}
    \par\noindent\textbf{\exlabel}~}
\newcommand{\exe}{\@ifstar{\@@exe}{\@exe}}

%    \end{macrocode}
% \end{macro}
% \begin{macro}{\points (ini)}
%    \begin{macrocode}    
\newcommand{\pointsstyle}{%
    \small\mdseries\sffamily\color{pointscolor}\fbox}
\newcommand*{\points}[1]{\hfill
    \pointsstyle{#1~%
        \ifthenelse{\lengthtest{#1 cm < 2cm}}{\pointname}{\pointsname}%
    }
}

%    \end{macrocode}
% Percent symbols are necessary to avoid spaces between the
% |\fbox| and its inner text.
% Note that, without |\lengthtest|, the test |#1 < 2| doesn't work with decimal numbers.
% % \end{macro}
%
% \subsection{Enumerations and lists}
%
% When using \textsf{babel} with the \texttt{french} option\footnote{The \texttt{french} option
% of \textsf{babel} supersedes the \texttt{frenchb} option.},
% \texttt{itemize} lists are modified with the same dash label for each list level. 
% These modifications are canceled here to restore default 
% \LaTeX\ \texttt{itemize} lists (labels and spaces).
% As for |\exetranslate|,
% we have create the |\standardfrenchlists| command who must be called into
% |AtBeginDocument| or not, depending on whether
% \textsf{exesheet} is loaded before \textsf{babel} or after.
%
% The |\setlist| command comes from the \textsf{enumitem} package
% (|\setenumerate| is deprecated). By default |itemsep=1ex| for lists of first level,
% and |leftmargin=1.5em| allows to align labels on the start of lines.
%
%    \begin{macrocode}
\newcommand\standardfrenchlists{%
    \@ifpackagewith{babel}{frenchb}{
        \frenchbsetup{StandardLists=true}}{}
    \@ifpackagewith{babel}{french}{
        \@ifundefined{frenchsetup}{
            \frenchbsetup{StandardLists=true}}{
            \frenchsetup{StandardLists=True}}
    }{}
}
\ifthenelse{\boolean{nosetlist}}{}{
    \AtBeginDocument{% if loaded before babel package
        \standardfrenchlists}
    \standardfrenchlists % necessary when loaded after babel
    \setlist[enumerate]{font=\bfseries}
    \setlist[enumerate,1]{topsep=1.5ex plus 1ex minus 1ex,leftmargin=1.5em}
}

%    \end{macrocode}
% \begin{macro}{\exenumerate}
%    \begin{macrocode}
\newenvironment{exenumerate}[1][]{%
    \setlist[enumerate]{font=\bfseries}
    \setlist[enumerate,1]{leftmargin=1.5em,
        itemsep=3ex plus 1ex minus 1ex,topsep=3ex plus 1ex minus 1ex}
    \setlist[enumerate,3]{noitemsep,nolistsep}
    \setlist[itemize]{noitemsep,nolistsep}
    \begin{enumerate}[#1]
        }{\end{enumerate}}

%    \end{macrocode}
% \end{macro}
% \DescribeEnv{tablenum1} \DescribeEnv{tablenuma}
% The |\NewTasks| command comes from the \textsf{tasks} package.
% It allows to define the environments \texttt{tablenum1}, \texttt{tablenuma}
% and \texttt{tablitem}.
% The horizontal spaces are adjusted to get a good alignment with
% items of other \texttt{enumerate} (or \texttt{itemize}) environments.
%
%    \begin{macrocode}
\ifthenelse{\boolean{nosetlist}}{
    \NewTasksEnvironment[label=\arabic*.,
        column-sep=1em,
        after-item-skip=0.5ex plus 0.5ex minus 0.5ex]{tablenum1}[\item](2)
    \NewTasksEnvironment[label=(\alph*),
        column-sep=1em,label-align=right,
        item-indent=2.15em,label-width=1.6em,label-offset=0.5em,
        after-item-skip=0.5ex plus 0.5ex minus 0.5ex]{tablenuma}[\item](2)
}{% by default
    \NewTasksEnvironment[label=\arabic*.,label-format=\bfseries,
        column-sep=1em,label-align=right,
        item-indent=1.5em,label-width=1em,label-offset=0.5em,
        after-item-skip=0.5ex plus 0.5ex minus 0.5ex]{tablenum1}[\item](2)
    \NewTasksEnvironment[label=(\alph*),label-format=\bfseries,
        column-sep=1em,label-align=right,
        item-indent=2.15em,label-width=1.6em,label-offset=0.5em,
        after-item-skip=0.5ex plus 0.5ex minus 0.5ex]{tablenuma}[\item](2)
}
\PackageWarning{exesheet}{Environment tablenum is deprecated and replaced by tablenum1}

%    \end{macrocode}
% \DescribeEnv{tablitem}
%    \begin{macrocode}
\NewTasksEnvironment[label=\labelitemi,
    label-align=right,
    item-indent=2.3333em,label-offset=0.5em,
    after-item-skip=0.5ex plus 0.5ex minus 0.5ex]{tablitem}[\item](2)

%    \end{macrocode}
% \DescribeEnv{colsenum} \DescribeEnv{colsenum*}
%    \begin{macrocode}
\newenvironment{colsenum}[2][]{%
    \setlength{\multicolsep}{2ex}
    \raggedcolumns % default is \flushcolumns
    \begin{multicols}{#2} % #2 = number of columns
        \begin{enumerate}[#1] % #1 = options of enumerate
            }{
        \end{enumerate}
    \end{multicols}
}

\newenvironment{colsenum*}[2][]{%
    \setlength{\multicolsep}{2ex}
    \begin{multicols}{#2} % #2 = number of columns
        \begin{enumerate}[#1] % #1 = options of enumerate
            }{
        \end{enumerate}
    \end{multicols}
}

%    \end{macrocode}
% \DescribeEnv{colsitem} \DescribeEnv{colsitem*}
%    \begin{macrocode}
\newenvironment{colsitem}[2][]{%
    \setlength{\multicolsep}{2ex}
    \raggedcolumns
    \begin{multicols}{#2}
        \begin{itemize}[#1]
            }{
        \end{itemize}
    \end{multicols}
}

\newenvironment{colsitem*}[2][]{%
    \setlength{\multicolsep}{2ex}
    \begin{multicols}{#2}
        \begin{itemize}[#1]
            }{
        \end{itemize}
    \end{multicols}
}

%    \end{macrocode}
%
% \subsection{Questions and answers}
%
% \DescribeMacro{\questionsonly} \DescribeMacro{\answersonly}
% The \texttt{questions} and \texttt{answers} booleans
% control the display of corresponding environments.
% When created, a boolean has the \texttt{false} value by default,
% but we define them to \texttt{true}.
% The |\questionsonly| and |\answersonly| macros are used as user interface
% to display only one of the two environments.
%    \begin{macrocode}
\newboolean{questions}
\newboolean{answers}
\setboolean{questions}{true}
\setboolean{answers}{true}
\newcommand{\questionsonly}{
    \setboolean{questions}{true}\setboolean{answers}{false}}
\newcommand{\answersonly}{
    \setboolean{questions}{false}\setboolean{answers}{true}}

%    \end{macrocode}
%
% \begin{environment}{questions}
% It is the |\comment| and |\endcomment| macros, provided by the  \textsf{versions} package,
% that allow the magic of conditional displays (we can also find them 
% in the \textsf{verbatim} or \textsf{version} packages).
% The noteworthy \textsf{codesection} package, allows to encapsulate 
% optional code between the macros |\BeginCodeSection|\marg{skip} 
% and |\EndCodeSection|\marg{skip},
% both in the text body and in the preamble,
% but these macros cannot be used inside an environment 
% as we did here for |\comment| and |\endcomment|.
%    \begin{macrocode}
\newcounter{exe@ini}
\newcounter{subpart@ini}

\newenvironment{questions}{
    \ifthenelse{\boolean{questions}}{
        \setcounter{exe@ini}{\value{exercise}}
        \setcounter{subpart@ini}{\value{subpart}}
    }{\comment}}%
{\ifthenelse{\boolean{questions}}{}{\endcomment}}

%    \end{macrocode}
% \end{environment}
% \begin{environment}{answers}
% The internal macro |\set@toclevel| calculates the title level
% of the word \emph{\correctionname} to display at the start of an
% \texttt{answers} environment (when \texttt{questions} and \texttt{answers} are
% displayed together). The principle is to compare the state of the counters
% \texttt{exercise} and \texttt{subpart} with those saved at the time of
% the call of \texttt{questions}. The |\@enumdepth| counter indicates the 
% \texttt{enumerate} list level in which we are (0 = out of lists).
% The optional parameter of the \texttt{answers} environment allows to force this title level.
%    \begin{macrocode}
\newcounter{@toclevel}
\newcommand{\set@toclevel}[1][]{
    \ifthenelse{\equal{#1}{}}{
        \ifthenelse{\value{exercise} > \value{exe@ini}}{
            \setcounter{@toclevel}{1}
        }{\ifthenelse{\equal{\the\@enumdepth}{0}}{
                % we're not in an enumerate environment
                \ifthenelse{\(\value{subpart} > \value{subpart@ini}\)
                    \or \(\value{subpart} = 0\)}{
                    \setcounter{@toclevel}{2}
                }{\setcounter{@toclevel}{3}}
            }{\setcounter{@toclevel}{4}}}
    }{\setcounter{@toclevel}{#1}}}

\newcommand{\correctionstyle}{\color{correctioncolor}}

\newenvironment{answers}[1][]{% #1 is the optional level
    \ifthenelse{\boolean{answers}}{%
        \ifthenelse{\boolean{questions}}{%
            \set@toclevel[#1]
            \ifthenelse{\value{@toclevel} = 1}{
                \section*{\correctionstyle\correctionname}
                \ifthenelse{\boolean{notoc}}{}{
                    \addcontentsline{toc}{section}{\correctionname}}
                \setcounter{exercise}{0}
            }{\ifthenelse{\value{@toclevel} = 2}{%
                    \subsection*{\correctionstyle\correctionname}
                    \ifthenelse{\boolean{notoc}}{}{
                        \addcontentsline{toc}{subsection}{\correctionname}}
                    \setcounter{subpart}{0}
                }{\ifthenelse{\value{@toclevel} = 3}{%
                        \subsubsection*{\correctionstyle\correctionname}
                        \ifthenelse{\boolean{notoc}}{}{
                            \addcontentsline{toc}{subsubsection}{
                                \correctionname}}
                    }{\par\textbf{\correctionstyle\correctionname}\par
                    }%
                }%
            }%
            \correctionstyle%
        }{}%
    }{\comment}
}{\ifthenelse{\boolean{answers}}{}{\endcomment}}

\newenvironment{answers*}{\ifthenelse{\boolean{answers}}{}{\comment}}%
{\ifthenelse{\boolean{answers}}{}{\endcomment}}

%    \end{macrocode}
% In the \texttt{answers} environment, if we place |\correctionstyle|
% before |\subsubsection|, the preceding vertical space may be too wide.
%
% \end{environment}
% \begin{macro}{\question}
%    \begin{macrocode}
\newcommand{\question}[2]{%
    \ifthenelse{\boolean{questions}}{#1}{}
    \bgroup
    \ifthenelse{\boolean{answers}}{
        \ifthenelse{\boolean{questions}}{
            \ifx#2\empty\else
                \par\correctionstyle\textbf{\correctionname}\par
            \fi
        }{}
        #2}{}
    \egroup
}

%    \end{macrocode}
% \end{macro}
% \begin{macro}{\points (patched)}
%    \begin{macrocode}
\let\@oldpoints\points
\renewcommand*{\points}[1]{%
    \ifthenelse{\boolean{questions}}{\@oldpoints{#1}}{}}

%    \end{macrocode}
% \end{macro}
%
% \subsection{Marginal notes}
%
% The commands |\displaypts|, |\displaypoints| and |\displaynotes|
% change the ratio between left and right margins\footnote{So 
% that the effect on the margin ratio is correct, these
% macros must be called, in the preamble, after other commands 
% that also could alter the page geometry.}.
%
% \DescribeMacro{\pts} \DescribeMacro{\displaypts}
%    \begin{macrocode}
\newboolean{marginpts}
\newcommand*{\pointmark}[1]{%
    \ifthenelse{\lengthtest{#1 cm < 2cm}}{#1 \ptname}{#1 \ptsname}}
\newcommand{\ptsstyle}[1]{%
    \footnotesize\centering\sffamily\color{ptscolor} (#1)}
\newcommand*{\pts}[1]{%
    \ifthenelse{\boolean{marginpts}}{%
        \mbox{}% 
        \marginpar{\hspace{0pt}%
            \ptsstyle{\pointmark{#1}}}%
    }{}%
    \ignorespaces
}
\newcommand{\displaypts}{%
    \reversemarginpar
    \geometry{hmarginratio=3:2}
    \setboolean{marginpts}{true}
}

%    \end{macrocode}
%
% \begin{macro}{\totalexe}
% In the following macros using
% |\marginpar|, percent symbols and |\ignorespaces| are necessary 
% to avoid too much space in the text (or the margin)
% where these macros are inserted.
%    \begin{macrocode}
\newlength{\ptsboxlength}
\setlength{\ptsboxlength}{3.1em}
\cornersize{1}
\newcommand*{\totalexe}[1]{% 
    \ifthenelse{\boolean{marginpoints}}{%
        \mbox{}%
        \marginpar{\markingstyle{\ovalbox{%
                    \makebox[\ptsboxlength]{\pointmark{#1}}}%
            }}%        
    }{}%
    \ignorespaces
}

%    \end{macrocode}
% \end{macro}
%
% \begin{macro}{\note}
% The booleans \texttt{marginpoints} and \texttt{marginfullnotes} control
% the display of marginal notes. If \texttt{marginpoints} is \texttt{false},
% \texttt{marginfullnotes} will be ignored.
% The |\noteragged| command is initialized outside |\displaynotes|
% such that we can use the |\note| command without |\displaynotes| (see further).
%    \begin{macrocode}
\newboolean{marginpoints}
\newboolean{marginfullnotes}

\newcommand{\markingstyle}[1]{\hspace{0pt}\footnotesize\sffamily%
    \centering\color{markingcolor}\textbf{#1}}
\newcommand{\noteragged}{\raggedleft}
\newcommand{\notestyle}[1]{\hspace{0pt}\footnotesize\sffamily%
    \noteragged\noindent\color{notecolor} #1}
\newcommand{\@note}[2][]{%
    \ifthenelse{\boolean{marginpoints}}{%
        \mbox{}%
        \marginpar{%
            \ifthenelse{\equal{#1}{}}{}{\markingstyle{#1}\\}%
            \ifthenelse{\boolean{marginfullnotes}}{\notestyle #2}{}%
        }%
    }{}%
    \ignorespaces
}
\newcommand{\@@note}[1]{%
    \ifthenelse{\boolean{marginpoints}}{%
        \mbox{}%
        \marginpar{\markingstyle{#1}}%
    }{}%
    \ignorespaces
}
\newcommand{\note}{\@ifstar{\@@note}{\@note}}

%    \end{macrocode}
% \end{macro}
%
% \DescribeMacro{\displaypoints} \DescribeMacro{\displaynotes}
%\DescribeMacro{\displaynotesright}
% In |\displaynotes|, the additional length \texttt{1\,in} matches 
% the default free space to the left of |\oddsidemargin|.
%
%    \begin{macrocode}
\newcommand{\displaypoints}{%
    \reversemarginpar
    \geometry{hmarginratio=3:2}
    \setboolean{marginpoints}{true}
}

\newcommand*{\displaynotes}[1][\raggedleft]{%
    \reversemarginpar
    \renewcommand{\noteragged}{#1}
    \geometry{hmarginratio=5:1}
    \setlength{\marginparwidth}{\oddsidemargin}
    \addtolength{\marginparwidth}{1in}
    \addtolength{\marginparwidth}{-\marginparsep}
    \setlength{\marginparwidth}{0.8\marginparwidth}
    \setboolean{marginpoints}{true}
    \setboolean{marginfullnotes}{true}
}

\newcommand*{\displaynotesright}[1][\raggedright]{%
    \normalmarginpar
    \renewcommand{\noteragged}{#1}
    \geometry{hmarginratio=1:5}
    \setlength{\marginparwidth}{\paperwidth}
    \addtolength{\marginparwidth}{-\textwidth}
    \addtolength{\marginparwidth}{-\oddsidemargin}
    \addtolength{\marginparwidth}{-\marginparsep}
    \addtolength{\marginparwidth}{-1in}
    \setlength{\marginparwidth}{0.8\marginparwidth}
    \setboolean{marginpoints}{true}
    \setboolean{marginfullnotes}{true}
}
%    \end{macrocode}
% The formatting of marginal notes can possibly be made more regular 
% with the \textsf{ragged2e} package, not loaded by \textsf{exesheet}.
% We will then have to enter, as an optional parameter of |\displaypoints|:
% |\RaggedLeft|, |\Centering|, |\RaggedRight| or |justifying|.
%
% \begin{macro}{\totalpoints}
%    \begin{macrocode}
\newcommand{\totalpoints}{%
    \ifthenelse{\boolean{marginpoints}}{\totalexe}{\points}}
%</package>
%    \end{macrocode}
% \end{macro}
%
% \Finale
\endinput
